\chapter[1903 --- Finsen]{1903: Niels Ryberg Finsen}
\label{ch:1903_nielsrybergfinsen}

\section*{Main Contribution}

\textit{Pioneered the use of concentrated light radiation to treat lupus vulgaris (skin tuberculosis), establishing phototherapy as a medical discipline.}

\section{Official Citation}

for his contribution to the treatment of diseases, especially lupus vulgaris

\section{Background and Historical Context}

Niels Ryberg Finsen was born on December 15, 1860, in T\`orshavn on the Faroe Islands, a remote North Atlantic archipelago under Danish sovereignty\index{Finsen, Niels Ryberg|textbf}. His early years on these windswept, often fog-shrouded islands, where sunlight was a precious and intermittent resource, were to have a significant influence on his life's work. The Faroe Islands experience was formative: Finsen grew up in a culture where sunlight was both celebrated and scarce, with the long winter nights fostering a deep appreciation for the healing properties of light. His father, a civil servant, emphasized education, and young Niels showed early promise in both science and literature.

Finsen's childhood in the Faroe Islands was marked by a distinctive relationship with light and darkness that would shape his scientific vision. The archipelago, located between Norway and Iceland at approximately 62°N latitude, experiences extreme photoperiod variations: in summer, the sun barely sets, providing nearly continuous twilight, while in winter, the islands endure months of profound darkness with only a few hours of dim daylight\index{Faroe Islands|textbf}. The local culture developed unique adaptations to this environment. Traditional Faroese homes were oriented to capture every available ray of light, and community life revolved around the seasons' light cycles. The annual \emph{óskafti} (wish season) in January, when the first glimpses of longer days return, was celebrated as a time of hope and renewal. These cultural patterns of light anticipation likely influenced Finsen's later sensitivity to light's therapeutic potential\index{Circadian rhythm|textbf}.

Finsen's father, Jóannes Finsen, served as a customs official and later as a government administrator. The family valued learning deeply, and Niels received an excellent education at the local school before moving to Copenhagen in 1877 at age 16 to attend the Latin School of the Cathedral of Our Lady\index{Education|textbf}. The transition from the Faroe Islands to the Danish capital was stark: Copenhagen, though more illuminated than T\`orshavn, still offered a completely different relationship with natural light. Finsen often reflected on how the Faroe Islands' unusual light conditions---both the midnight sun and the polar night---had sensitized him to phenomena that city-dwelling scientists overlooked. He carried with him a deep awareness of light as both life-giving force and healing element, an insight that would later inform his scientific inquiries into the therapeutic potential of concentrated radiation\index{Light therapy|textbf}.

Finsen studied at the University of Copenhagen, initially focusing on mathematics and philosophy before completing his medical degree in 1890. During his studies, he became fascinated by the work of Thomas Hunt Morgan and other naturalists who studied the effects of light on biological organisms. His early research interests encompassed optics, the physics of radiation, and the relationship between light exposure and human health. After graduating, Finsen worked briefly at the University Hospital in Copenhagen before establishing his own private practice, where he began to explore the therapeutic potential of light systematically\index{University of Copenhagen|textbf}.

The medical context of Finsen's work was shaped by two converging developments in the late nineteenth century. First, the bacteriological revolution led by Koch and Pasteur had established that many diseases were caused by specific microorganisms, and that these organisms could be killed or inhibited in the laboratory\index{Robert Koch|textbf}. Second, the physics of light and radiation was advancing rapidly. The electromagnetic spectrum had been characterized, X-rays had been discovered by Wilhelm Röntgen in 1895, and researchers were beginning to explore the biological effects of different wavelengths of light.

The bactericidal properties of sunlight had been observed as early as the 1870s, when the French physician Armand Dubois demonstrated that sunlight could kill bacteria, and Otto Friedrich Mebes showed that light could inhibit the growth of microorganisms\index{Bactericidal properties|textbf}. However, these observations had not been translated into therapeutic applications. Finsen recognized that if sunlight could kill bacteria, concentrated light might be able to cure diseases caused by bacteria.

Lupus vulgaris, the form of cutaneous tuberculosis that was the primary target of Finsen's therapy, was a disfiguring and chronic disease that affected thousands of people in Europe\index{Lupus vulgaris|textbf}. Caused by \emph{Mycobacterium tuberculosis}, the same organism that caused pulmonary tuberculosis, lupus vulgaris produced progressive destruction of skin and underlying tissues, particularly on the face. The disease was notoriously difficult to treat; surgical excision was often ineffective, and the available chemical treatments were painful and of limited efficacy. Patients with lupus vulgaris were often severely disfigured, and social stigma compounded the physical suffering.

\section{The Discovery}

Finsen's path to phototherapy began with observations about the effects of light on the skin\index{Phototherapy|textbf}. He noticed that people living in dark, cramped conditions were more susceptible to certain skin diseases and that sunlight appeared to have therapeutic effects. Through careful experimentation, Finsen systematically investigated which wavelengths of light were most effective against bacteria and diseased tissue.

Finsen's key experimental breakthrough was the demonstration that ultraviolet (UV) radiation, particularly in the wavelength range of 300--400 nanometers, possessed powerful bactericidal properties. He showed that UV light could kill the tubercle bacillus and other pathogenic organisms in vitro, and that the effect was dependent on both the intensity and the wavelength of the light\index{Ultraviolet radiation|textbf}. Red and yellow light had little bactericidal effect, while blue and especially ultraviolet light were highly effective.

However, translating these laboratory findings into a practical therapy posed significant technical challenges. Sunlight was variable and unreliable, particularly in Scandinavian winters. The UV component of sunlight was attenuated by glass windows, and the light intensity at the skin surface was insufficient to produce a therapeutic effect in a reasonable treatment time\index{Technical challenges|textbf}. Finsen addressed these challenges through a series of ingenious engineering solutions.

Finsen's technical innovations were remarkable for their ingenuity and systematic approach\index{Technical innovations|textbf}. Working initially in a modest laboratory converted from his home, Finsen developed three principal instruments. First, he created the Finsen cap, a device for irradiating the scalp and face that used a system of mirrors and lenses to concentrate light onto the affected area while protecting the patient's eyes. Second, he designed the Finsen lamp, a high-intensity light source that produced concentrated ultraviolet radiation. Finsen experimented with various light sources, including sunlight, electric arc lamps, and mercury vapor lamps. He ultimately developed a system that used a powerful carbon arc lamp as the light source, with lenses and filters to concentrate the UV radiation and eliminate the heat-producing infrared wavelengths that could damage the skin\index{Finsen lamp|textbf}. The system included quartz condensing lenses, which transmitted UV radiation efficiently (unlike ordinary glass, which blocked UV), and water filters to absorb heat.

Third, Finsen invented the contact lens technique, or \emph{linspresse}, a method for increasing the penetration of UV radiation into the skin. He observed that the pressure of the lens against the skin blanched the tissue, displacing blood from the superficial capillaries. Since blood absorbs UV radiation, this blanching allowed the light to penetrate more deeply into the diseased tissue\index{Photodynamic effect|textbf}. This seemingly simple observation---that hemoglobin in blood competed with the therapeutic effect of UV light---was a critical insight that dramatically improved the efficacy of phototherapy. The technique required the patient to remain perfectly still for extended periods, often several hours per session.

Finsen also pioneered the use of standardized dosimetry in phototherapy. He developed a system for measuring and recording the intensity and duration of light exposure, allowing treatments to be reproduced and optimized. This quantitative approach was revolutionary in a field that had previously relied on qualitative judgments about light exposure. His careful documentation of treatment parameters established the foundation for evidence-based phototherapy.

The clinical application of Finsen's phototherapy was painstaking and labor-intensive. Treatment sessions lasted several hours, as the light had to be applied to small areas of diseased tissue sequentially\index{Clinical application|textbf}. Patients underwent multiple sessions over weeks or months. The results, however, were remarkable. Finsen demonstrated that concentrated UV radiation could cure lupus vulgaris---a disease that had previously been virtually untreatable. He published his clinical results beginning in 1896, reporting cure rates of over 80 percent in patients with lupus vulgaris. His report on the first 80 patients treated at the Finsen Institute in Copenhagen, published in 1901, showed that the majority of patients were cured or significantly improved.

Finsen's clinical work required extraordinary patience and precision. Each treatment session could last up to three hours, during which the patient had to remain completely still while the concentrated light was applied to the affected area. The equipment required constant adjustment, and Finsen personally supervised each procedure. Despite the demanding nature of the treatment, patients reported significant improvement, and many achieved complete clearance of their lesions\index{Patient care|textbf}.

Finsen also explored the therapeutic applications of UV light beyond lupus vulgaris. He investigated its effects on other skin diseases, including scrofuloderma (cutaneous tuberculosis), certain forms of eczema, and alopecia areata. He also studied the general physiological effects of light exposure, including its influence on blood composition and immune function---areas that would later be developed by other researchers into the field of photobiology\index{Photobiology|textbf}.

\section{Impact on Modern Medicine}

Finsen's phototherapy opened an entirely new avenue of medical treatment---the therapeutic use of light---that has expanded dramatically in the century since his Nobel Prize-winning work\index{Light therapy|textbf}. The most direct descendant of Finsen's work is ultraviolet phototherapy, which remains a mainstay of treatment for several skin diseases. Narrowband UVB phototherapy, which uses light in the 311--313 nm wavelength range---remarkably close to the wavelengths that Finsen identified as most bactericidal---is now the standard treatment for psoriasis, eczema (atopic dermatitis), vitiligo, and several other dermatological conditions\index{Psoriasis|textbf}. The treatment is administered in controlled doses using specialized UV lamps, and it is both safe and effective when properly administered.

Photodynamic therapy (PDT), a modern treatment for certain cancers and precancerous conditions, builds on Finsen's principle that light can be used to destroy diseased tissue\index{Photodynamic therapy|textbf}. In PDT, a photosensitizing agent is administered to the patient and accumulates preferentially in tumor cells. When the tumor is then exposed to light of a specific wavelength, the photosensitizer absorbs the light energy and generates reactive oxygen species that destroy the cancer cells. PDT is used to treat certain forms of skin cancer, esophageal cancer, lung cancer, and bladder cancer.

Beyond dermatology and oncology, light-based therapies have found applications in many areas of medicine. Laser therapy is used in ophthalmology, surgery, and dermatology\index{Laser therapy|textbf}. Photobiomodulation therapy (low-level laser therapy) is used for wound healing, pain management, and inflammation. Blue light therapy is used to treat neonatal jaundice by converting bilirubin into a water-soluble form that can be excreted by the infant\index{Neonatal jaundice|textbf}. In each of these applications, the fundamental principle is the same one that Finsen articulated: that specific wavelengths of light can produce specific biological effects.

Finsen's work also contributed to the broader understanding of the biological effects of light, including the discovery of vitamin D synthesis in the skin by ultraviolet radiation---a finding that has significant implications for bone health, immune function, and the prevention of rickets\index{Vitamin D|textbf}. The discovery that UV light triggers vitamin D production in 1922 by Harry Steenbock built directly on Finsen's foundational work on UV radiation and established light's role in nutritional medicine.

\section{Legacy}

Niels Ryberg Finsen received the Nobel Prize in Physiology or Medicine in 1903, at the age of 42, making him one of the youngest recipients of the award\index{Nobel Prize|textbf}. He was also one of the first scientists from the Nordic countries to receive a Nobel Prize, and his achievement brought international attention to the scientific community of Denmark and Scandinavia.

Finsen's career was cut short by illness. He suffered from a chronic cardiovascular condition that progressively weakened him throughout his working years. He died on September 24, 1904, at the age of 43, just over a year after receiving the Nobel Prize. Despite his short life, his impact on medicine was significant and enduring\index{Cardiovascular disease|textbf}.

The Finsen Institute in Copenhagen, which he founded in 1896, continued to be a center for phototherapy research and clinical treatment for many decades\index{Finsen Institute|textbf}. It eventually became part of the larger Copenhagen University Hospital system, where its legacy of light-based medical research persists.

Finsen's story also illustrates an important principle in the history of medicine: that the translation of basic scientific observations into practical therapies often requires not only scientific insight but also engineering ingenuity and clinical determination\index{Translational research|textbf}. Finsen's understanding of the bactericidal properties of UV light was not, in itself, sufficient to cure lupus vulgaris. It was his engineering innovations---the concentrated light source, the quartz lenses, the water filters, the pressure technique---that made phototherapy a practical reality. This combination of scientific curiosity, technical skill, and clinical vision is a model for translational research that remains relevant today\index{Engineering|textbf}.

The Finsen Institute that he founded in Copenhagen in 1896 became a model for specialized disease treatment centers\index{Finsen Institute|textbf}. It was the first hospital in the world dedicated exclusively to the treatment of a single disease using a specific physical modality. The institute's success demonstrated the value of specialized research hospitals, a model that has been widely adopted in modern medicine\index{Medical research|textbf}.

Finsen's influence extended beyond his immediate clinical work. He trained numerous physicians from around the world who came to study at the Finsen Institute, spreading the technique globally. His published works, including \emph{Collected Papers} (published posthumously in 1906), became essential references for researchers working in photobiology\index{Publications|textbf}.

\section{Modern Developments}

Phototherapy has advanced far beyond Finsen's original arc lamp\index{LED therapy|textbf}. The development of light-emitting diodes (LEDs) in the late twentieth century revolutionized phototherapy, replacing bulky arc lamps with compact, efficient, wavelength-specific sources. Modern narrowband UVB (311 nm) LED systems provide precise dosing with minimal side effects, and are the standard of care for psoriasis and vitiligo treatment worldwide. Photodynamic therapy (PDT) uses light-activated drugs to destroy cancer cells and is approved for esophageal, lung, and bladder cancers. Low-level laser therapy treats chronic wounds and inflammatory conditions. In 2020, optogenetics---using light to control neurons---enabled researchers to restore partial vision in blind mice, opening new frontiers in neuroscience\index{Optogenetics|textbf}. UV-C disinfection technology, based on the same light-tissue interactions Finsen studied, became critical during the COVID-19 pandemic\index{UV-C disinfection|textbf}. The modern Finsen Memorial in Copenhagen, inaugurated in 1933, stands as a testament to his enduring legacy as the father of phototherapy\index{Finsen Memorial|textbf}.

Today, phototherapy is estimated to benefit millions of patients worldwide annually\index{Global health|textbf}. The International Commission on Illumination (CIE) has established standardized protocols for phototherapy based on Finsen's foundational principles. Research continues to expand the applications of light-based treatments, including studies into seasonal affective disorder (SAD), psoriasis, and certain types of cancer. Finsen's insight that light could be a therapeutic tool---rather than merely a diagnostic aid---has generated a medical discipline that continues to evolve and save lives more than a century after his Nobel Prize.


\section{Bibliography}

\begin{thebibliography}{9}
\bibitem{nobel-1903}
Nobel Prize Outreach, ``The Nobel Prize in Physiology or Medicine 1903,''\emph{NobelPrize.org}.\\
\url{https://www.nobelprize.org/prizes/medicine/1903/summary/}

\bibitem{lecture-1903}
Niels Ryberg Finsen, ``Nobel Lecture,''\emph{NobelPrize.org}. Winner-authored primary source; downloadable from the official Nobel Prize archive.\\
\url{https://www.nobelprize.org/prizes/medicine/1903/finsen/lecture/}
\end{thebibliography}
